\documentclass[11pt,a4paper]{article}
\usepackage[utf8]{inputenc}
\usepackage[T1]{fontenc}
\usepackage[spanish,es-minimal,es-noshorthands]{babel}
\usepackage{lmodern}
\usepackage[sfdefault,lining]{FiraSans}
\usepackage[a4paper,top=2.3cm,bottom=2.2cm,left=1.8cm,right=1.8cm,headheight=1.3cm,headsep=0.4cm,footskip=1cm]{geometry}
\usepackage{amsmath,amssymb}
\allowdisplaybreaks
\usepackage{icomma}
\usepackage{graphicx}
\usepackage[table]{xcolor}
\usepackage{tikz}
\usepackage{fontawesome5}
\usepackage{ccicons}
\usepackage{fancyhdr}
\usepackage{lastpage}
\usepackage{tasks}
\usepackage{needspace}
\usepackage[most]{tcolorbox}
\usepackage{array}
\usepackage[hidelinks]{hyperref}
\AtBeginDocument{\hypersetup{pdftitle={Ficha de ejercicios de la unidad 1: Matrices y determinantes. 2º Bach.},pdfsubject={Matemáticas CCSS 2º Bach.},pdfkeywords={matrices, operaciones con matrices, potencias de matrices, determinantes, regla de Sarrus, adjuntos, propiedades de los determinantes, rango, matriz inversa, ecuaciones matriciales, sistemas matriciales}}}

\definecolor{navy}{HTML}{1B2748}
\definecolor{accent}{HTML}{E8603A}
\definecolor{teal}{HTML}{1F9E90}
\definecolor{soft}{HTML}{EEF1F8}
\definecolor{muted}{HTML}{5D6685}

\setlength{\parindent}{0pt}
\setlength{\parskip}{3pt}
\settasks{label=\alph*), label-format=\bfseries\color{navy}, label-width=1.3em, item-indent=1.8em, column-sep=0.6em, after-item-skip=5pt, before-skip=3pt, after-skip=4pt}

% ---------- capçalera i peu ----------
\newcommand{\llicencia}{\href{https://creativecommons.org/licenses/by-nc-sa/4.0/deed.es}{\ccbyncsa}\ \ Licencia Creative Commons Atribución-NoComercial-CompartirIgual 4.0 Internacional (CC BY-NC-SA 4.0)}
\pagestyle{fancy}
\fancyhf{}
\fancyhead[L]{\raisebox{-0.35\height}{\includegraphics[height=0.95cm]{logo-gva.png}}}
\fancyhead[C]{\small\color{muted}Ficha de ejercicios · Unidad 1: Matrices y determinantes · 2\textsuperscript{o} Bach.}
\fancyhead[R]{\raisebox{-0.35\height}{\includegraphics[height=0.95cm]{logo-vermellar.png}}}
\fancyfoot[L]{\scriptsize\color{muted}\llicencia}
\fancyfoot[R]{\small\color{muted}\thepage}
\renewcommand{\headrulewidth}{0.4pt}
\renewcommand{\headrule}{\color{muted!40}\hrule height \headrulewidth}
\fancypagestyle{primera}{\fancyhf{}\renewcommand{\headrulewidth}{0pt}
  \fancyfoot[L]{\scriptsize\color{muted}\llicencia}
  \fancyfoot[R]{\small\color{muted}\thepage}}

% ---------- nivells i exercicis ----------
\newcommand{\estrelles}[1]{{\color{accent}\ifcase#1\or\faStar\faStar[regular]\faStar[regular]\or\faStar\faStar\faStar[regular]\or\faStar\faStar\faStar\fi}}
\newcounter{exer}
\newcommand{\ex}[1]{\par\needspace{4\baselineskip}\medskip\refstepcounter{exer}\noindent{\color{navy}\bfseries\theexer.}\ \ }

\newcommand{\seccio}[2]{\par\needspace{10\baselineskip}\bigskip
  \begin{tcolorbox}[colback=navy,colframe=navy,coltext=white,arc=2mm,boxrule=0pt,left=3mm,right=3mm,top=1.2mm,bottom=1.2mm]
  {\large\bfseries #1\quad #2}
  \end{tcolorbox}\vspace{-2pt}}


% quadre de recordatori (només a la fitxa d'exercicis)
\newcommand{\recorda}[1]{\par\smallskip\begin{tcolorbox}[colback=accent!7,colframe=accent!55,arc=2mm,boxrule=0.5pt,left=2.5mm,right=2.5mm,top=1mm,bottom=1mm]\small{\color{accent}\faLightbulb[regular]}\ \textbf{Recuerda:}\ #1\end{tcolorbox}}
\newcommand{\res}[1]{#1}\newcommand{\rest}[1]{#1}

% =====================================================================
% Capçalera FentMates de les fitxes: un color per a cada curs.
% Els colors són els de la pàgina inicial (LEVELS a assets/js/materials.js);
% el logo ve de _eines/logo/pdf/ (python3 _eines/logo/genera.py).
%
% Ús, al final del preàmbul de la unitat:
%   % =====================================================================
% Capçalera FentMates de les fitxes: un color per a cada curs.
% Els colors són els de la pàgina inicial (LEVELS a assets/js/materials.js);
% el logo ve de _eines/logo/pdf/ (python3 _eines/logo/genera.py).
%
% Ús, al final del preàmbul de la unitat:
%   % =====================================================================
% Capçalera FentMates de les fitxes: un color per a cada curs.
% Els colors són els de la pàgina inicial (LEVELS a assets/js/materials.js);
% el logo ve de _eines/logo/pdf/ (python3 _eines/logo/genera.py).
%
% Ús, al final del preàmbul de la unitat:
%   \input{../comu/fentmates.tex}
%   \fmcurs{1eso}
% i, al començament del document:
%   \capcalera{Matemàtiques}{1r ESO}{Fitxa d'exercicis · Unitat 2}{Divisibilitat}
% Opcional abans de \input: \def\fmneed{14} (espai mínim abans de cada secció).
% =====================================================================
\makeatletter
\def\fm@logodir{../../logo/pdf}
\providecommand{\fmneed}{10}
\def\fm@cursnom{}

% clau del curs, color i color del text damunt del color
\newcommand{\fm@curs}[3]{\@namedef{fm@color@#1}{#2}\@namedef{fm@text@#1}{#3}}
\fm@curs{1eso}{E6B84A}{navy}
\fm@curs{2eso}{3B82F6}{white}
\fm@curs{3eso}{FF7A50}{white}
\fm@curs{4eso-a}{A78BFA}{white}
\fm@curs{4eso-b}{6D28D9}{white}
\fm@curs{1batx-mi}{9B1C1C}{white}
\fm@curs{1batx-ccss}{EF6B6B}{white}
\fm@curs{2batx-mii}{1E6B47}{white}
\fm@curs{2batx-ccss}{2BB3A3}{white}

\newcommand{\fmcurs}[1]{%
  \@ifundefined{fm@color@#1}{\PackageError{fentmates}{Curs desconegut: #1}{}}{}%
  \edef\fm@tmp{\noexpand\definecolor{curs}{HTML}{\@nameuse{fm@color@#1}}%
               \noexpand\colorlet{curstext}{\@nameuse{fm@text@#1}}}\fm@tmp
  \colorlet{cursfosc}{curs!72!black}%
  \def\fm@logo{\fm@logodir/vertical-#1.pdf}%
  \def\fm@marca{\fm@logodir/marca-#1.pdf}}

% ---------- capçalera de la primera pàgina ----------
% #1 matèria (petit), #2 curs, #3 tipus de document i unitat, #4 títol
\newcommand{\capcalera}[4]{%
  \gdef\fm@cursnom{#2}%
  \par\noindent
  \begin{tcolorbox}[colback=curs,colframe=curs,boxrule=0pt,arc=2.2mm,width=3.5cm,height=2.15cm,
      valign=center,left=3mm,right=2mm,top=1mm,bottom=1mm,nobeforeafter,box align=center]
    \color{curstext}{\footnotesize #1}\par\vspace{1pt}{\fontsize{19}{21}\selectfont\bfseries #2\par}%
  \end{tcolorbox}\hspace{5mm}%
  \begin{minipage}[c]{\dimexpr\linewidth-3.5cm-5mm-2.9cm\relax}
    {\color{cursfosc}\small\bfseries\MakeUppercase{#3}}\par\vspace{2pt}
    {\color{navy}\fontsize{20}{23}\selectfont\bfseries #4\par}
  \end{minipage}\hfill
  \begin{minipage}[c]{2.6cm}\raggedleft\includegraphics[width=2.6cm]{\fm@logo}\end{minipage}\par
  \vspace{2pt}}

% ---------- barres de secció del color del curs ----------
\renewcommand{\seccio}[2]{\par\needspace{\fmneed\baselineskip}\bigskip
  \begin{tcolorbox}[colback=curs,colframe=curs,coltext=curstext,arc=2mm,boxrule=0pt,left=3mm,right=3mm,top=1.2mm,bottom=1.2mm]
  {\large\bfseries #1\quad #2}
  \end{tcolorbox}\vspace{-2pt}}

% ---------- estrelles de dificultat del color del curs (a Batxillerat no se'n posen) ----------
\@ifundefined{estrelles}{}{%
  \renewcommand{\estrelles}[1]{{\color{curs}\ifcase#1\or\faStar\faStar[regular]\faStar[regular]\or\faStar\faStar\faStar[regular]\or\faStar\faStar\faStar\fi}}}

% ---------- capçalera i peu de les pàgines ----------
% primera pàgina: logos institucionals al peu
\fancypagestyle{primera}{\fancyhf{}\renewcommand{\headrulewidth}{0pt}
  \fancyfoot[L]{\raisebox{-0.32\height}{\includegraphics[height=0.78cm]{logo-gva.png}}\hspace{5mm}%
    \raisebox{-0.32\height}{\includegraphics[height=0.78cm]{logo-vermellar.png}}}
  \fancyfoot[R]{\scriptsize\color{muted}\href{https://creativecommons.org/licenses/by-nc-sa/4.0/deed.ca}{\ccbyncsa}\ \ CC BY-NC-SA 4.0 · fentmates.org\qquad\small\thepage}}
% la resta: el símbol a l'esquerra i el curs a la dreta, amb el títol al centre com abans
\fancyhead[L]{\raisebox{-0.3\height}{\includegraphics[height=0.62cm]{\fm@marca}}}
\fancyhead[R]{\small\bfseries\color{cursfosc}\fm@cursnom}
\renewcommand{\headrule}{\color{curs!45}\hrule height \headrulewidth}
\makeatother

%   \fmcurs{1eso}
% i, al començament del document:
%   \capcalera{Matemàtiques}{1r ESO}{Fitxa d'exercicis · Unitat 2}{Divisibilitat}
% Opcional abans de \input: \def\fmneed{14} (espai mínim abans de cada secció).
% =====================================================================
\makeatletter
\def\fm@logodir{../../logo/pdf}
\providecommand{\fmneed}{10}
\def\fm@cursnom{}

% clau del curs, color i color del text damunt del color
\newcommand{\fm@curs}[3]{\@namedef{fm@color@#1}{#2}\@namedef{fm@text@#1}{#3}}
\fm@curs{1eso}{E6B84A}{navy}
\fm@curs{2eso}{3B82F6}{white}
\fm@curs{3eso}{FF7A50}{white}
\fm@curs{4eso-a}{A78BFA}{white}
\fm@curs{4eso-b}{6D28D9}{white}
\fm@curs{1batx-mi}{9B1C1C}{white}
\fm@curs{1batx-ccss}{EF6B6B}{white}
\fm@curs{2batx-mii}{1E6B47}{white}
\fm@curs{2batx-ccss}{2BB3A3}{white}

\newcommand{\fmcurs}[1]{%
  \@ifundefined{fm@color@#1}{\PackageError{fentmates}{Curs desconegut: #1}{}}{}%
  \edef\fm@tmp{\noexpand\definecolor{curs}{HTML}{\@nameuse{fm@color@#1}}%
               \noexpand\colorlet{curstext}{\@nameuse{fm@text@#1}}}\fm@tmp
  \colorlet{cursfosc}{curs!72!black}%
  \def\fm@logo{\fm@logodir/vertical-#1.pdf}%
  \def\fm@marca{\fm@logodir/marca-#1.pdf}}

% ---------- capçalera de la primera pàgina ----------
% #1 matèria (petit), #2 curs, #3 tipus de document i unitat, #4 títol
\newcommand{\capcalera}[4]{%
  \gdef\fm@cursnom{#2}%
  \par\noindent
  \begin{tcolorbox}[colback=curs,colframe=curs,boxrule=0pt,arc=2.2mm,width=3.5cm,height=2.15cm,
      valign=center,left=3mm,right=2mm,top=1mm,bottom=1mm,nobeforeafter,box align=center]
    \color{curstext}{\footnotesize #1}\par\vspace{1pt}{\fontsize{19}{21}\selectfont\bfseries #2\par}%
  \end{tcolorbox}\hspace{5mm}%
  \begin{minipage}[c]{\dimexpr\linewidth-3.5cm-5mm-2.9cm\relax}
    {\color{cursfosc}\small\bfseries\MakeUppercase{#3}}\par\vspace{2pt}
    {\color{navy}\fontsize{20}{23}\selectfont\bfseries #4\par}
  \end{minipage}\hfill
  \begin{minipage}[c]{2.6cm}\raggedleft\includegraphics[width=2.6cm]{\fm@logo}\end{minipage}\par
  \vspace{2pt}}

% ---------- barres de secció del color del curs ----------
\renewcommand{\seccio}[2]{\par\needspace{\fmneed\baselineskip}\bigskip
  \begin{tcolorbox}[colback=curs,colframe=curs,coltext=curstext,arc=2mm,boxrule=0pt,left=3mm,right=3mm,top=1.2mm,bottom=1.2mm]
  {\large\bfseries #1\quad #2}
  \end{tcolorbox}\vspace{-2pt}}

% ---------- estrelles de dificultat del color del curs (a Batxillerat no se'n posen) ----------
\@ifundefined{estrelles}{}{%
  \renewcommand{\estrelles}[1]{{\color{curs}\ifcase#1\or\faStar\faStar[regular]\faStar[regular]\or\faStar\faStar\faStar[regular]\or\faStar\faStar\faStar\fi}}}

% ---------- capçalera i peu de les pàgines ----------
% primera pàgina: logos institucionals al peu
\fancypagestyle{primera}{\fancyhf{}\renewcommand{\headrulewidth}{0pt}
  \fancyfoot[L]{\raisebox{-0.32\height}{\includegraphics[height=0.78cm]{logo-gva.png}}\hspace{5mm}%
    \raisebox{-0.32\height}{\includegraphics[height=0.78cm]{logo-vermellar.png}}}
  \fancyfoot[R]{\scriptsize\color{muted}\href{https://creativecommons.org/licenses/by-nc-sa/4.0/deed.ca}{\ccbyncsa}\ \ CC BY-NC-SA 4.0 · fentmates.org\qquad\small\thepage}}
% la resta: el símbol a l'esquerra i el curs a la dreta, amb el títol al centre com abans
\fancyhead[L]{\raisebox{-0.3\height}{\includegraphics[height=0.62cm]{\fm@marca}}}
\fancyhead[R]{\small\bfseries\color{cursfosc}\fm@cursnom}
\renewcommand{\headrule}{\color{curs!45}\hrule height \headrulewidth}
\makeatother

%   \fmcurs{1eso}
% i, al començament del document:
%   \capcalera{Matemàtiques}{1r ESO}{Fitxa d'exercicis · Unitat 2}{Divisibilitat}
% Opcional abans de \input: \def\fmneed{14} (espai mínim abans de cada secció).
% =====================================================================
\makeatletter
\def\fm@logodir{../../logo/pdf}
\providecommand{\fmneed}{10}
\def\fm@cursnom{}

% clau del curs, color i color del text damunt del color
\newcommand{\fm@curs}[3]{\@namedef{fm@color@#1}{#2}\@namedef{fm@text@#1}{#3}}
\fm@curs{1eso}{E6B84A}{navy}
\fm@curs{2eso}{3B82F6}{white}
\fm@curs{3eso}{FF7A50}{white}
\fm@curs{4eso-a}{A78BFA}{white}
\fm@curs{4eso-b}{6D28D9}{white}
\fm@curs{1batx-mi}{9B1C1C}{white}
\fm@curs{1batx-ccss}{EF6B6B}{white}
\fm@curs{2batx-mii}{1E6B47}{white}
\fm@curs{2batx-ccss}{2BB3A3}{white}

\newcommand{\fmcurs}[1]{%
  \@ifundefined{fm@color@#1}{\PackageError{fentmates}{Curs desconegut: #1}{}}{}%
  \edef\fm@tmp{\noexpand\definecolor{curs}{HTML}{\@nameuse{fm@color@#1}}%
               \noexpand\colorlet{curstext}{\@nameuse{fm@text@#1}}}\fm@tmp
  \colorlet{cursfosc}{curs!72!black}%
  \def\fm@logo{\fm@logodir/vertical-#1.pdf}%
  \def\fm@marca{\fm@logodir/marca-#1.pdf}}

% ---------- capçalera de la primera pàgina ----------
% #1 matèria (petit), #2 curs, #3 tipus de document i unitat, #4 títol
\newcommand{\capcalera}[4]{%
  \gdef\fm@cursnom{#2}%
  \par\noindent
  \begin{tcolorbox}[colback=curs,colframe=curs,boxrule=0pt,arc=2.2mm,width=3.5cm,height=2.15cm,
      valign=center,left=3mm,right=2mm,top=1mm,bottom=1mm,nobeforeafter,box align=center]
    \color{curstext}{\footnotesize #1}\par\vspace{1pt}{\fontsize{19}{21}\selectfont\bfseries #2\par}%
  \end{tcolorbox}\hspace{5mm}%
  \begin{minipage}[c]{\dimexpr\linewidth-3.5cm-5mm-2.9cm\relax}
    {\color{cursfosc}\small\bfseries\MakeUppercase{#3}}\par\vspace{2pt}
    {\color{navy}\fontsize{20}{23}\selectfont\bfseries #4\par}
  \end{minipage}\hfill
  \begin{minipage}[c]{2.6cm}\raggedleft\includegraphics[width=2.6cm]{\fm@logo}\end{minipage}\par
  \vspace{2pt}}

% ---------- barres de secció del color del curs ----------
\renewcommand{\seccio}[2]{\par\needspace{\fmneed\baselineskip}\bigskip
  \begin{tcolorbox}[colback=curs,colframe=curs,coltext=curstext,arc=2mm,boxrule=0pt,left=3mm,right=3mm,top=1.2mm,bottom=1.2mm]
  {\large\bfseries #1\quad #2}
  \end{tcolorbox}\vspace{-2pt}}

% ---------- estrelles de dificultat del color del curs (a Batxillerat no se'n posen) ----------
\@ifundefined{estrelles}{}{%
  \renewcommand{\estrelles}[1]{{\color{curs}\ifcase#1\or\faStar\faStar[regular]\faStar[regular]\or\faStar\faStar\faStar[regular]\or\faStar\faStar\faStar\fi}}}

% ---------- capçalera i peu de les pàgines ----------
% primera pàgina: logos institucionals al peu
\fancypagestyle{primera}{\fancyhf{}\renewcommand{\headrulewidth}{0pt}
  \fancyfoot[L]{\raisebox{-0.32\height}{\includegraphics[height=0.78cm]{logo-gva.png}}\hspace{5mm}%
    \raisebox{-0.32\height}{\includegraphics[height=0.78cm]{logo-vermellar.png}}}
  \fancyfoot[R]{\scriptsize\color{muted}\href{https://creativecommons.org/licenses/by-nc-sa/4.0/deed.ca}{\ccbyncsa}\ \ CC BY-NC-SA 4.0 · fentmates.org\qquad\small\thepage}}
% la resta: el símbol a l'esquerra i el curs a la dreta, amb el títol al centre com abans
\fancyhead[L]{\raisebox{-0.3\height}{\includegraphics[height=0.62cm]{\fm@marca}}}
\fancyhead[R]{\small\bfseries\color{cursfosc}\fm@cursnom}
\renewcommand{\headrule}{\color{curs!45}\hrule height \headrulewidth}
\makeatother

\fmcurs{2batx-ccss}

\begin{document}
\thispagestyle{primera}

\capcalera{Matemáticas CCSS}{2\textsuperscript{o} Bach.}{Ficha de ejercicios · Unidad 1}{Matrices y determinantes}
\vspace{10pt}
Nombre y apellidos: \rule{6.6cm}{0.4pt}\hfill Grupo: \rule{1.6cm}{0.4pt}\hfill Fecha: \rule{2.2cm}{0.4pt}
\vspace{8pt}

\begin{tcolorbox}[colback=soft,colframe=soft,arc=2mm,left=3mm,right=3mm,top=1.5mm,bottom=1.5mm]
\small Resuelve los ejercicios en el cuaderno mostrando todos los pasos. Antes de multiplicar dos matrices, comprueba que los órdenes lo permiten, y recuerda que el producto de matrices \textbf{no} es conmutativo: en las ecuaciones matriciales, fíjate bien por qué lado multiplicas por la inversa.
\end{tcolorbox}

% =====================================================
\seccio{1}{Matrices. Definiciones y tipos}

\ex{1} Escribe las matrices siguientes:\begin{tasks}(1)
  \task $A=(a_{ij})$, de orden $3\times4$, con $a_{ij}=2i-j$. ¿Cuál es el elemento $a_{23}$? ¿Y el elemento $a_{32}$?
  \task $B=(b_{ij})$, cuadrada de orden 3, con $b_{ij}=i+j$ si $i\leq j$ \ y \ $b_{ij}=0$ si $i>j$. ¿De qué tipo es?
  \task $C=(c_{ij})$, de orden $2\times3$, con $c_{ij}=(-1)^{i+j}\cdot i$.
\end{tasks}\recorda{$a_{ij}$ es el elemento que ocupa la fila $i$ y la columna $j$. Una matriz de orden $n\times m$ tiene $n$ filas y $m$ columnas.}

\ex{1} Indica de qué tipo es cada una de las matrices siguientes (fila, columna, cuadrada, nula, diagonal, escalar, triangular, simétrica, antisimétrica):\begin{tasks}(4)
  \task $\begin{pmatrix}1 & -2 & 5\end{pmatrix}$
  \task $\begin{pmatrix}3\\0\\-1\end{pmatrix}$
  \task $\begin{pmatrix}2 & 0 & 0\\0 & -1 & 0\\0 & 0 & 4\end{pmatrix}$
  \task $\begin{pmatrix}3 & 0\\0 & 3\end{pmatrix}$
  \task $\begin{pmatrix}1 & 4 & -2\\0 & 3 & 5\\0 & 0 & -1\end{pmatrix}$
  \task $\begin{pmatrix}2 & -1 & 3\\-1 & 0 & 5\\3 & 5 & 4\end{pmatrix}$
  \task $\begin{pmatrix}0 & 2 & -3\\-2 & 0 & 1\\3 & -1 & 0\end{pmatrix}$
  \task $\begin{pmatrix}0 & 0\\0 & 0\end{pmatrix}$
\end{tasks}

\ex{1} Calcula $x$, $y$, $z$ y $t$ para que se cumpla la igualdad \ $\begin{pmatrix}x+y & 2\\ 3 & x-y\end{pmatrix}=\begin{pmatrix}5 & z-1\\ 2t+1 & 1\end{pmatrix}$.

\ex{2} Halla los valores de las letras para que:\begin{tasks}(1)
  \task $A=\begin{pmatrix}1 & a+1 & 3\\ 2 & 0 & b\\ c-1 & 5 & 4\end{pmatrix}$ \ sea simétrica.
  \task $B=\begin{pmatrix}0 & x & -2\\ 3 & y & z\\ 2 & -1 & 0\end{pmatrix}$ \ sea antisimétrica.
\end{tasks}

\ex{1} Una cadena tiene tres tiendas de bicicletas, B1, B2 y B3. En enero, B1 vendió 12 bicicletas urbanas y 8 de montaña; B2, 5 urbanas y 10 de montaña, y B3, 9 urbanas y 4 de montaña.\begin{tasks}(1)
  \task Escribe la matriz $V$ de ventas (una fila por tienda y una columna por tipo). ¿Cuál es su orden?
  \task ¿Qué significa el elemento $v_{32}$?
  \task Escribe $V^t$. ¿Qué información da cada fila de $V^t$?
\end{tasks}


% =====================================================
\seccio{2}{Operaciones con matrices}

\ex{1} Dadas las matrices \ $A=\begin{pmatrix}2 & -1 & 0\\1 & 3 & -2\end{pmatrix}$, \ $B=\begin{pmatrix}1 & 0 & 4\\-2 & 1 & 3\end{pmatrix}$ \ y \ $C=\begin{pmatrix}0 & 2\\-1 & 1\end{pmatrix}$, \ calcula:\begin{tasks}(3)
  \task $A+B$
  \task $2A-3B$
  \task $A^t-B^t$
  \task $C\cdot A$
  \task $A\cdot B^t$
  \task $C^2-I$
\end{tasks}\recorda{Para sumar, las matrices deben tener el mismo orden. El producto $A\cdot B$ solo existe si el número de \textbf{columnas} de $A$ es igual al número de \textbf{filas} de $B$: \ $(n\times m)\cdot(m\times p)=n\times p$.}

\ex{1} Las matrices $A$, $B$, $C$ y $D$ tienen órdenes $2\times3$, $3\times3$, $3\times2$ y $2\times2$, respectivamente. Indica cuáles de estos productos se pueden hacer y, si se pueden hacer, qué orden tiene el resultado:\begin{tasks}(4)
  \task $A\cdot B$
  \task $B\cdot A$
  \task $A\cdot C$
  \task $C\cdot A$
  \task $B\cdot C$
  \task $C\cdot D$
  \task $D\cdot A$
  \task $A^2$
\end{tasks}

\ex{1} Sean \ $A=\begin{pmatrix}1 & 2\\3 & -1\end{pmatrix}$ \ y \ $B=\begin{pmatrix}0 & 1\\2 & -2\end{pmatrix}$.\begin{tasks}(1)
  \task Calcula $A\cdot B$ y $B\cdot A$. ¿Es conmutativo el producto de matrices?
  \task Calcula $(A+B)^2$ y $A^2+2AB+B^2$. ¿Por qué no coinciden?
\end{tasks}

\ex{1} Calcula $A\cdot B$ y $B\cdot A$ para \ $A=\begin{pmatrix}1 & 0 & 2\\-1 & 1 & 0\\0 & 3 & 1\end{pmatrix}$ \ y \ $B=\begin{pmatrix}2 & 1 & 0\\0 & -1 & 1\\1 & 0 & 3\end{pmatrix}$.

\ex{1} Sean la matriz fila $F=\begin{pmatrix}1 & 2 & -1\end{pmatrix}$ y la matriz columna $C=\begin{pmatrix}3\\0\\2\end{pmatrix}$. Calcula $F\cdot C$ y $C\cdot F$.

\ex{2} Dada \ $A=\begin{pmatrix}1 & 1\\0 & 1\end{pmatrix}$, calcula $A^2$, $A^3$ y $A^4$. Deduce la expresión de $A^n$ y calcula $A^{20}$.\recorda{Para calcular potencias altas, calcula unas cuantas, busca la regularidad y generalízala.}

\ex{2} Dada \ $B=\begin{pmatrix}0 & 1 & 0\\0 & 0 & 1\\1 & 0 & 0\end{pmatrix}$, calcula $B^2$ y $B^3$. ¿Cuánto vale $B^{100}$?

\ex{2} Sea \ $C=\begin{pmatrix}0 & 1 & 2\\0 & 0 & 3\\0 & 0 & 0\end{pmatrix}$. Calcula $C^2$ y $C^3$. ¿Cuánto vale $C^n$ para $n\geq3$? Calcula también $(I+C)^2$.

\ex{3} Halla todas las matrices $X$ de orden 2 que conmutan con \ $A=\begin{pmatrix}1 & 1\\0 & 1\end{pmatrix}$, es decir, que cumplen $A\cdot X=X\cdot A$.

\ex{2} Calcula $a$ y $b$ para que la matriz \ $A=\begin{pmatrix}a & 1\\ 0 & b\end{pmatrix}$ \ cumpla \ $A^2=\begin{pmatrix}4 & 3\\ 0 & 1\end{pmatrix}$.

\ex{2} Comprueba que la matriz \ $A=\begin{pmatrix}2 & -1\\3 & -2\end{pmatrix}$ \ cumple $A^2=I$ y calcula $A^{51}$ y $A^{100}$.

\ex{2} Una matriz cuadrada es \emph{idempotente} si $A^2=A$. Comprueba que \ $A=\begin{pmatrix}2 & -2 & -4\\-1 & 3 & 4\\1 & -2 & -3\end{pmatrix}$ \ es idempotente y calcula $A^n$.


% =====================================================
\seccio{3}{Determinantes}

\ex{1} Calcula los determinantes de orden 2 siguientes:\begin{tasks}(4)
  \task $\begin{vmatrix}3 & -2\\5 & 4\end{vmatrix}$
  \task $\begin{vmatrix}-1 & 7\\2 & -3\end{vmatrix}$
  \task $\begin{vmatrix}0 & 5\\-6 & 9\end{vmatrix}$
  \task $\begin{vmatrix}x & x+1\\ x-1 & x\end{vmatrix}$
\end{tasks}\recorda{$\begin{vmatrix}a & b\\ c & d\end{vmatrix}=ad-bc$. \ Solo tienen determinante las matrices cuadradas.}

\ex{1} Calcula, con la regla de Sarrus:\begin{tasks}(3)
  \task $\begin{vmatrix}1 & 2 & 3\\0 & -1 & 4\\2 & 1 & 0\end{vmatrix}$
  \task $\begin{vmatrix}2 & -1 & 1\\3 & 0 & 2\\1 & 4 & -1\end{vmatrix}$
  \task $\begin{vmatrix}1 & 1 & 1\\2 & 3 & 4\\4 & 9 & 16\end{vmatrix}$
\end{tasks}\recorda{Sarrus: sumamos los productos de la diagonal principal y de las dos ``paralelas'' y restamos los de la diagonal secundaria y sus ``paralelas''.}

\ex{2} Dada la matriz \ $A=\begin{pmatrix}2 & -1 & 3\\0 & 4 & 0\\1 & 5 & -2\end{pmatrix}$:\begin{tasks}(1)
  \task Calcula los adjuntos $A_{12}$, $A_{22}$ y $A_{31}$.
  \task Calcula $|A|$ desarrollando por la fila que te resulte más cómoda.
\end{tasks}\recorda{El adjunto del elemento $a_{ij}$ es $A_{ij}=(-1)^{i+j}\cdot\alpha_{ij}$, donde $\alpha_{ij}$ es el menor que se obtiene al suprimir la fila $i$ y la columna $j$.}

\ex{3} Calcula el determinante de orden 4 \ $\begin{vmatrix}2 & 0 & 1 & 0\\1 & 3 & 0 & 0\\0 & 1 & 2 & 1\\1 & 0 & 1 & 0\end{vmatrix}$ \ desarrollando por la columna que tenga más ceros.

\ex{2} Resuelve las ecuaciones siguientes:\begin{tasks}(2)
  \task $\begin{vmatrix}x & 3\\ 2 & x+1\end{vmatrix}=0$
  \task $\begin{vmatrix}x & 2 & 1\\ 1 & x & 0\\ 0 & 1 & 1\end{vmatrix}=3$
\end{tasks}

\ex{2} Estudia para qué valores de $m$ se anulan los determinantes siguientes:\begin{tasks}(2)
  \task $\begin{vmatrix}m+1 & 3\\ 1 & m-1\end{vmatrix}$
  \task $\begin{vmatrix}1 & 2 & m\\ 2 & m & 1\\ 1 & 1 & 1\end{vmatrix}$
\end{tasks}


% =====================================================
\seccio{4}{Propiedades de los determinantes}

\ex{2} Sabiendo que \ $\begin{vmatrix}a & b & c\\ d & e & f\\ g & h & i\end{vmatrix}=5$, \ calcula, sin desarrollarlos:\begin{tasks}(2)
  \task $\begin{vmatrix}d & e & f\\ a & b & c\\ g & h & i\end{vmatrix}$
  \task $\begin{vmatrix}a & b & c\\ 2d & 2e & 2f\\ g & h & i\end{vmatrix}$
  \task $\begin{vmatrix}a & d & g\\ b & e & h\\ c & f & i\end{vmatrix}$
  \task $\begin{vmatrix}a & b & c\\ d+3a & e+3b & f+3c\\ g & h & i\end{vmatrix}$
  \task $\begin{vmatrix}-a & -b & -c\\ -d & -e & -f\\ -g & -h & -i\end{vmatrix}$
  \task $\begin{vmatrix}2a & 2b & 2c\\ 2d & 2e & 2f\\ 2g & 2h & 2i\end{vmatrix}$
  \task $\begin{vmatrix}a & b & c\\ d & e & f\\ 2g-a & 2h-b & 2i-c\end{vmatrix}$
\end{tasks}\recorda{Cambiar dos filas de sitio cambia el signo; multiplicar una fila por $k$ multiplica el determinante por $k$; sumar a una fila una combinación lineal de las otras no lo cambia; $|A^t|=|A|$; \ $|k\cdot A_n|=k^n|A|$.}

\ex{2} Sean $A$ y $B$ dos matrices cuadradas de orden 3 con $|A|=2$ y $|B|=-3$. Calcula:\begin{tasks}(3)
  \task $|A\cdot B|$
  \task $|3A|$
  \task $|A^{-1}|$
  \task $|A^t\cdot B|$
  \task $|A^2\cdot B^{-1}|$
  \task $|-2B|$
\end{tasks}\recorda{$|A\cdot B|=|A|\cdot|B|$, \ \ $|A^{-1}|=\dfrac{1}{|A|}$, \ \ $|A^t|=|A|$, \ \ $|k\cdot A_n|=k^n|A|$.}

\ex{2} Justifica, sin desarrollarlos, que los determinantes siguientes valen cero:\begin{tasks}(3)
  \task $\begin{vmatrix}1 & 2 & 3\\ 4 & 5 & 6\\ 7 & 8 & 9\end{vmatrix}$
  \task $\begin{vmatrix}2 & -1 & 4\\ 6 & -3 & 12\\ 1 & 5 & 0\end{vmatrix}$
  \task $\begin{vmatrix}1 & a & b+c\\ 1 & b & a+c\\ 1 & c & a+b\end{vmatrix}$
\end{tasks}

\ex{3} Calcula aplicando las propiedades (haz ceros antes de desarrollar):\begin{tasks}(2)
  \task $\begin{vmatrix}1 & 1 & 1 & 1\\1 & 2 & 2 & 2\\1 & 2 & 3 & 3\\1 & 2 & 3 & 4\end{vmatrix}$
  \task $\begin{vmatrix}x & 1 & 1\\ 1 & x & 1\\ 1 & 1 & x\end{vmatrix}$ \ (factoriza el resultado)
\end{tasks}

\ex{3} Razona las afirmaciones siguientes:\begin{tasks}(1)
  \task Si $A$ es una matriz antisimétrica de orden 3, entonces $|A|=0$.
  \task Si $A$ es cuadrada y $A\cdot A^t=I$, entonces $|A|=1$ o $|A|=-1$.
  \task Si $A$ es idempotente ($A^2=A$), entonces $|A|$ solo puede valer 0 o 1.
\end{tasks}


% =====================================================
\seccio{5}{Rango de una matriz}

\ex{1} Calcula por el método de Gauss el rango de las matrices siguientes:\begin{tasks}(3)
  \task $\begin{pmatrix}2 & 4 & 6\\1 & 2 & 3\\3 & 6 & 9\end{pmatrix}$
  \task $\begin{pmatrix}1 & -1 & 2 & 0\\2 & 1 & 0 & 3\\3 & 0 & 2 & 3\end{pmatrix}$
  \task $\begin{pmatrix}1 & 2 & 1\\0 & 1 & 3\\2 & 1 & 0\\1 & 0 & 2\end{pmatrix}$
\end{tasks}\recorda{El rango es el número de filas (o columnas) linealmente independientes. Por Gauss: escalonamos la matriz y contamos las filas que no son nulas.}

\ex{2} Calcula el rango de \ $A=\begin{pmatrix}1 & 2 & 0 & 1\\2 & 1 & 3 & -1\\3 & 3 & 3 & 0\end{pmatrix}$ \ utilizando menores (determinantes).\recorda{El rango es el orden del mayor menor no nulo. Si un menor de orden $k$ es no nulo, basta con \emph{orlarlo} (añadirle una fila y una columna) para buscar uno de orden $k+1$.}

\ex{3} Estudia el rango de la matriz \ $A=\begin{pmatrix}1 & 1 & k\\ 1 & k & 1\\ k & 1 & 1\end{pmatrix}$ \ según los valores de $k$.

\ex{3} Estudia el rango de la matriz \ $B=\begin{pmatrix}1 & 2 & -1 & 3\\ 2 & 1 & a & 1\\ 1 & -1 & 3 & -2\end{pmatrix}$ \ según los valores de $a$.


% =====================================================
\seccio{6}{Matriz inversa}

\ex{1} Calcula, a partir de la definición ($A\cdot A^{-1}=I$, planteando un sistema), la inversa de las matrices siguientes, si existe:\begin{tasks}(2)
  \task $A=\begin{pmatrix}2 & 1\\5 & 3\end{pmatrix}$
  \task $B=\begin{pmatrix}1 & 2\\2 & 4\end{pmatrix}$
\end{tasks}\recorda{Una matriz cuadrada tiene inversa si y solo si su determinante es distinto de cero (matriz \emph{regular}).}

\ex{1} Calcula la inversa de las matrices siguientes por el método de los adjuntos, si existe, y comprueba el resultado:\begin{tasks}(3)
  \task $\begin{pmatrix}4 & 3\\2 & 2\end{pmatrix}$
  \task $\begin{pmatrix}-1 & 2\\3 & -5\end{pmatrix}$
  \task $\begin{pmatrix}6 & 4\\3 & 2\end{pmatrix}$
\end{tasks}\recorda{$A^{-1}=\dfrac{1}{|A|}\big(\mathrm{Adj}(A)\big)^t$. \ Para orden 2: \ $\begin{pmatrix}a & b\\ c & d\end{pmatrix}^{-1}=\dfrac{1}{ad-bc}\begin{pmatrix}d & -b\\ -c & a\end{pmatrix}$.}

\ex{2} \label{ex:inv3}Calcula la inversa de las matrices siguientes:\begin{tasks}(2)
  \task $A=\begin{pmatrix}1 & 2 & 1\\0 & 1 & 0\\2 & 0 & 3\end{pmatrix}$
  \task $B=\begin{pmatrix}2 & 1 & 0\\0 & 1 & 3\\1 & 0 & 1\end{pmatrix}$
\end{tasks}

\ex{2} Dada la matriz \ $M=\begin{pmatrix}1 & 0 & -1\\ 0 & k & 3\\ 4 & 1 & -k\end{pmatrix}$:\begin{tasks}(1)
  \task ¿Para qué valores de $k$ tiene inversa?
  \task Calcula $M^{-1}$ para $k=2$.
\end{tasks}

\ex{3} Sea \ $A=\begin{pmatrix}2 & 1\\0 & 1\end{pmatrix}$.\begin{tasks}(1)
  \task Comprueba que $A^2-3A+2I=O$.
  \task Utiliza esta igualdad para calcular $A^{-1}$ sin determinantes.
\end{tasks}

\ex{2} Dadas \ $A=\begin{pmatrix}1 & 2\\1 & 3\end{pmatrix}$ \ y \ $B=\begin{pmatrix}2 & 1\\1 & 1\end{pmatrix}$, calcula $(A\cdot B)^{-1}$ y comprueba que coincide con $B^{-1}\cdot A^{-1}$, pero no con $A^{-1}\cdot B^{-1}$.


% =====================================================
\seccio{7}{Ecuaciones y sistemas matriciales}

\ex{2} Dadas \ $A=\begin{pmatrix}2 & 1\\5 & 3\end{pmatrix}$ \ y \ $B=\begin{pmatrix}1 & 0\\2 & -1\end{pmatrix}$, resuelve las ecuaciones matriciales:\begin{tasks}(2)
  \task $A\cdot X=B$
  \task $X\cdot A=B$
\end{tasks}\recorda{Despejamos $X$ igual que en una ecuación numérica, pero \textbf{no podemos dividir}: multiplicamos por la inversa por el mismo lado donde está la matriz. \ $AX=B\Rightarrow X=A^{-1}B$; \ \ $XA=B\Rightarrow X=BA^{-1}$.}

\ex{2} Resuelve la ecuación $A\cdot X+B=C$, siendo \ $A=\begin{pmatrix}1 & 2\\1 & 3\end{pmatrix}$, \ $B=\begin{pmatrix}1 & 0\\2 & 1\end{pmatrix}$ \ y \ $C=\begin{pmatrix}3 & 4\\5 & 9\end{pmatrix}$.

\ex{3} Resuelve la ecuación $X\cdot A-2X=B$, siendo \ $A=\begin{pmatrix}4 & 1\\3 & 3\end{pmatrix}$ \ y \ $B=\begin{pmatrix}1 & 2\\0 & -1\end{pmatrix}$.

\ex{3} Resuelve la ecuación $A\cdot X\cdot B=C$, siendo \ $A=\begin{pmatrix}1 & 1\\0 & 1\end{pmatrix}$, \ $B=\begin{pmatrix}2 & 1\\1 & 1\end{pmatrix}$ \ y \ $C=\begin{pmatrix}3 & 2\\1 & 1\end{pmatrix}$.

\ex{2} Halla la matriz $X$ que cumple $A\cdot X=B$, siendo \ $A=\begin{pmatrix}1 & 2 & 1\\0 & 1 & 0\\2 & 0 & 3\end{pmatrix}$ \ y \ $B=\begin{pmatrix}1 & 0\\2 & 1\\3 & -1\end{pmatrix}$. \ (La inversa de $A$ la tienes en el ejercicio \ref{ex:inv3}.)

\ex{2} Resuelve el sistema matricial \ $\left\{\begin{aligned}2X+Y&=\begin{pmatrix}3 & 1\\0 & 2\end{pmatrix}\\ X-Y&=\begin{pmatrix}0 & -1\\-3 & 4\end{pmatrix}\end{aligned}\right.$\recorda{Los sistemas matriciales se resuelven por reducción o sustitución, como los numéricos.}

\ex{2} Resuelve el sistema matricial \ $\left\{\begin{aligned}3X-2Y&=\begin{pmatrix}4 & -5\\-6 & 3\end{pmatrix}\\ X+Y&=\begin{pmatrix}3 & 0\\3 & 1\end{pmatrix}\end{aligned}\right.$

\ex{3} Dada \ $A=\begin{pmatrix}1 & -1\\2 & -1\end{pmatrix}$, halla la matriz $X$ tal que \ $A\cdot X=A^t+I$.


% =====================================================
\seccio{8}{Problemas de aplicación}

\ex{2} Una empresa fabrica tres modelos de muebles (A, B y C) en dos fábricas. La matriz $P$ da las unidades fabricadas cada semana (filas: fábricas F1 y F2; columnas: modelos A, B y C) y la matriz $C$, las horas de trabajo y el coste de material (en euros) de cada unidad (filas: modelos A, B y C):\[P=\begin{pmatrix}40 & 25 & 10\\30 & 35 & 20\end{pmatrix}\qquad C=\begin{pmatrix}2 & 15\\3 & 20\\5 & 30\end{pmatrix}\]\begin{tasks}(1)
  \task Calcula $P\cdot C$ y explica qué significa cada elemento.
  \task ¿Tiene sentido calcular $C\cdot P$?
\end{tasks}\recorda{Para multiplicar, las columnas de la primera matriz y las filas de la segunda deben referirse a la \textbf{misma} cosa (aquí, los modelos).}

\ex{2} Tres tiendas (B1, B2 y B3) venden pantalones, camisas y chaquetas. Las matrices $V$ y $W$ recogen las unidades vendidas en enero y en febrero (filas: tiendas; columnas: pantalones, camisas, chaquetas), y $p$ es la matriz columna de los precios: 30 €, 20 € y 50 €.\[V=\begin{pmatrix}10 & 15 & 4\\8 & 20 & 6\\12 & 10 & 5\end{pmatrix}\qquad W=\begin{pmatrix}12 & 10 & 6\\9 & 18 & 4\\10 & 14 & 7\end{pmatrix}\qquad p=\begin{pmatrix}30\\20\\50\end{pmatrix}\]\begin{tasks}(1)
  \task Calcula $V+W$. ¿Qué representa?
  \task Calcula $(V+W)\cdot p$ e interprétalo.
  \task Si en marzo todos los precios suben un 10\,\%, ¿qué matriz da los nuevos precios?
\end{tasks}

\ex{3} Una compañía aérea une Valencia (V), Madrid (M), Barcelona (B) y Palma (P) con estos vuelos directos: de V a M, B y P; de M a V y B; de B a M y P, y de P a V.\begin{tasks}(1)
  \task Escribe la matriz $A$ en la que $a_{ij}=1$ si hay vuelo directo de la ciudad $i$ a la ciudad $j$, y $a_{ij}=0$ si no lo hay (orden: V, M, B, P).
  \task Calcula $A^2$. ¿Qué indican sus elementos?
  \task ¿De cuántas maneras se puede ir de Palma a Barcelona haciendo exactamente una escala? ¿Y salir de Valencia y volver haciendo una escala?
\end{tasks}

\ex{3} Una comarca tiene dos zonas, la costa (C) y el interior (I). Cada año, el 10\,\% de los habitantes de la costa se traslada al interior, y el 20\,\% de los del interior se traslada a la costa; el resto se queda donde está. Ahora hay 5000 habitantes en la costa y 3000 en el interior.\begin{tasks}(1)
  \task Escribe la matriz $T$ tal que $\begin{pmatrix}C_{n+1}\\ I_{n+1}\end{pmatrix}=T\cdot\begin{pmatrix}C_n\\ I_n\end{pmatrix}$.
  \task ¿Cuántos habitantes habrá en cada zona dentro de un año? ¿Y dentro de dos años?
  \task Comprueba que, si hubiera $\tfrac{16000}{3}$ habitantes en la costa y $\tfrac{8000}{3}$ en el interior, la distribución no cambiaría.
\end{tasks}

\vfill
\begin{center}\color{muted}\small\faCheckCircle[regular]\ \ ¡Ánimo! Comprueba tus resultados con el laboratorio de matrices de los apuntes interactivos.\end{center}

\end{document}
